% OncoRx-Bench -- CAFCW26/SC26 review manuscript
% SC26 requires the IEEE conference proceedings format.
% Review mode is anonymous by default. Change \anonymoustrue to
% \anonymousfalse only for a camera-ready version after confirming the
% CAFCW26 review policy.
% CAFCW26 proceedings papers must meet SC26 technical-paper requirements.
% Prepare the mandatory paper checklist and Artifact Description appendix
% (or an explanation that no artifact exists). Confirm only the workshop-
% specific upload mechanism and deadlines in Linklings.
\documentclass[conference]{IEEEtran}

\newif\ifanonymous
\anonymoustrue

\usepackage{array}
\usepackage{booktabs}
\usepackage{cite}
\usepackage{enumitem}
\usepackage{graphicx}
\usepackage[nopatch=eqnum]{microtype}
\usepackage{tabularx}
\usepackage{tikz}
\usepackage{xcolor}
\usepackage{url}
\usetikzlibrary{arrows.meta,positioning}
\usepackage[hidelinks]{hyperref}

\newcolumntype{Y}{>{\raggedright\arraybackslash}X}
\newcommand{\dataset}{\textsc{OncoRx-Bench}}
\newcommand{\code}[1]{\texttt{#1}}
\ifanonymous
\newcommand{\kbfull}{an anonymized integrated oncology knowledge artifact}
\newcommand{\kbname}{the anonymized upstream knowledge artifact}
\newcommand{\kbadj}{upstream-artifact}
\else
\newcommand{\kbfull}{the Unified Oncology Treatment Database (UOTD)}
\newcommand{\kbname}{UOTD}
\newcommand{\kbadj}{UOTD}
\fi
\IfFileExists{generated_stats.tex}
  {% Generated by scripts/profile_release.py; do not edit by hand.
\newcommand{\FinalDrugObjects}{3,525}
\newcommand{\FinalUniqueDrugs}{170}
\newcommand{\FinalUniqueEvidenceSurfaces}{433}
\newcommand{\FinalUniqueExplicitSurfaces}{328}
\newcommand{\FinalUniqueInferenceSurfaces}{105}
\newcommand{\FinalSigObjects}{1,286}
\newcommand{\FinalRegimenObjects}{444}
\newcommand{\FinalUniqueRegimens}{176}
\newcommand{\FinalComponentAssignments}{1,110}
\newcommand{\FinalUniqueComponents}{109}
\newcommand{\FinalMeanTerms}{10.1}
\newcommand{\FinalMedianTerms}{8.0}
\newcommand{\FinalDrugViewRows}{241}
\newcommand{\FinalRegimenViewRows}{492}
\newcommand{\FinalConditionViewRows}{804}
\newcommand{\FinalCOneMeanTerms}{6.6}
\newcommand{\FinalCOneDrugs}{770}
\newcommand{\FinalCOneSigs}{354}
\newcommand{\FinalCOneRegimens}{0}
\newcommand{\FinalCTwoMeanTerms}{13.2}
\newcommand{\FinalCTwoDrugs}{530}
\newcommand{\FinalCTwoSigs}{515}
\newcommand{\FinalCTwoRegimens}{0}
\newcommand{\FinalCThreeMeanTerms}{12.4}
\newcommand{\FinalCThreeDrugs}{1,410}
\newcommand{\FinalCThreeSigs}{267}
\newcommand{\FinalCThreeRegimens}{394}
\newcommand{\FinalCFourMeanTerms}{8.1}
\newcommand{\FinalCFourDrugs}{460}
\newcommand{\FinalCFourSigs}{12}
\newcommand{\FinalCFourRegimens}{22}
\newcommand{\FinalCFiveMeanTerms}{11.2}
\newcommand{\FinalCFiveDrugs}{355}
\newcommand{\FinalCFiveSigs}{138}
\newcommand{\FinalCFiveRegimens}{28}
\newcommand{\FinalValidationErrors}{0}
\newcommand{\FinalDuplicateTexts}{0}
\newcommand{\FinalSplitLeakageErrors}{0}
\newcommand{\FinalJsonlSha}{97a6c2a22eacf0d68e1762051fbe57b943f24c032a014eadd64f7e36aecf952f}
\newcommand{\FinalJsonlShaA}{97a6c2a22eacf0d68e1762051fbe57b9}
\newcommand{\FinalJsonlShaB}{43f24c032a014eadd64f7e36aecf952f}
}
  {% Generated by scripts/profile_release.py; do not edit by hand.
\newcommand{\FinalDrugObjects}{3,525}
\newcommand{\FinalUniqueDrugs}{170}
\newcommand{\FinalUniqueEvidenceSurfaces}{433}
\newcommand{\FinalUniqueExplicitSurfaces}{328}
\newcommand{\FinalUniqueInferenceSurfaces}{105}
\newcommand{\FinalSigObjects}{1,286}
\newcommand{\FinalRegimenObjects}{444}
\newcommand{\FinalUniqueRegimens}{176}
\newcommand{\FinalComponentAssignments}{1,110}
\newcommand{\FinalUniqueComponents}{109}
\newcommand{\FinalMeanTerms}{10.1}
\newcommand{\FinalMedianTerms}{8.0}
\newcommand{\FinalDrugViewRows}{241}
\newcommand{\FinalRegimenViewRows}{492}
\newcommand{\FinalConditionViewRows}{804}
\newcommand{\FinalCOneMeanTerms}{6.6}
\newcommand{\FinalCOneDrugs}{770}
\newcommand{\FinalCOneSigs}{354}
\newcommand{\FinalCOneRegimens}{0}
\newcommand{\FinalCTwoMeanTerms}{13.2}
\newcommand{\FinalCTwoDrugs}{530}
\newcommand{\FinalCTwoSigs}{515}
\newcommand{\FinalCTwoRegimens}{0}
\newcommand{\FinalCThreeMeanTerms}{12.4}
\newcommand{\FinalCThreeDrugs}{1,410}
\newcommand{\FinalCThreeSigs}{267}
\newcommand{\FinalCThreeRegimens}{394}
\newcommand{\FinalCFourMeanTerms}{8.1}
\newcommand{\FinalCFourDrugs}{460}
\newcommand{\FinalCFourSigs}{12}
\newcommand{\FinalCFourRegimens}{22}
\newcommand{\FinalCFiveMeanTerms}{11.2}
\newcommand{\FinalCFiveDrugs}{355}
\newcommand{\FinalCFiveSigs}{138}
\newcommand{\FinalCFiveRegimens}{28}
\newcommand{\FinalValidationErrors}{0}
\newcommand{\FinalDuplicateTexts}{0}
\newcommand{\FinalSplitLeakageErrors}{0}
\newcommand{\FinalJsonlSha}{97a6c2a22eacf0d68e1762051fbe57b943f24c032a014eadd64f7e36aecf952f}
\newcommand{\FinalJsonlShaA}{97a6c2a22eacf0d68e1762051fbe57b9}
\newcommand{\FinalJsonlShaB}{43f24c032a014eadd64f7e36aecf952f}
}
\newcommand{\blfootnote}[1]{%
  \begingroup
  \renewcommand{\thefootnote}{}%
  \footnote{#1}%
  \addtocounter{footnote}{-1}%
  \endgroup
}

\title{OncoRx-Bench: A Knowledge-Grounded Framework for
Generating Synthetic Benchmarks for Oncology Clinical Entity Extraction}

\ifanonymous
\author{\IEEEauthorblockN{Anonymous Authors}
\IEEEauthorblockA{Anonymous Institution}}
\else
\author{\IEEEauthorblockN{Abhishek Shivanna, Jordan Tschida, Mayanka Chandrashekar,\\
John Gounley, and Heidi A. Hanson}
\IEEEauthorblockA{Computational Sciences and Engineering Division\\
Oak Ridge National Laboratory}}
\fi

\begin{document}
\maketitle

\begin{abstract}
Expert-labeled clinical corpora are costly to create and difficult to share,
yet evaluating large language models (LLMs) benefits from repeatable benchmarks
that expose specific extraction challenges. Oncology is especially demanding
because a
system must recognize and normalize drug variants, recover prescribing
attributes, distinguish contextual mentions, and resolve named regimens into
their component agents. We present \dataset, a reusable knowledge-grounded
framework for constructing controlled synthetic benchmarks for these tasks.
The framework separates a versioned knowledge-input layer from a configurable
challenge taxonomy, template library, constrained samplers, scenario-specific
text and annotation generation, integrity controls, and
export interfaces. The reference knowledge views are materialized by a
checksum-verifying adapter from a pinned revision of \kbfull{}. A reference
configuration uses 347 templates across five categories and 20 subcategories to
generate 2,000 short clinical-style snippets. The resulting dataset contains
\FinalDrugObjects{} structured drug objects spanning
\FinalUniqueDrugs{} normalized drugs and
\FinalUniqueEvidenceSurfaces{} stored evidence surfaces
(\FinalUniqueExplicitSurfaces{} literal drug surfaces and
\FinalUniqueInferenceSurfaces{} regimen-evidence spans), together with
\FinalRegimenObjects{} regimen objects covering
\FinalUniqueRegimens{} regimen names. Its
task-selected targets are intended to support future benchmarking of LLM-based
entity recognition, normalization, prescribing-attribute extraction,
contextual status interpretation, regimen resolution, and robustness to
abbreviations, brands, and misspellings. Within each scenario-specific
function, synthetic text and task-selected target objects are produced from a
shared pool of sampled values, reducing the amount of sample-by-sample manual
annotation needed for new configurations. Synthetic benchmarks remain
complementary to
expert-labeled clinical data and require independent clinical review before any
claim of clinical validity. This paper describes the framework and benchmark
interfaces; systematic LLM comparisons are reserved for future work.
\end{abstract}

\section{Introduction}
\label{sec:introduction}

Oncology notes compress medication information into generic and brand names,
abbreviations, regimen acronyms, dose--route--frequency instructions, cycle
descriptions, prior therapy, temporary holds, negation, allergies, and adverse
reactions. Medication extraction benchmarks have therefore progressed from
drug-mention recognition toward attribute and context modeling. The i2b2
medication challenge included dosage, route, frequency, duration, and reason
\cite{uzuner2010}; MADE 1.0 added indications and adverse drug events in
cancer-patient records \cite{jagannatha2019}; and the 2022 n2c2 task modeled
medication-event action, negation, temporality, certainty, and actor
\cite{mahajan2023}. These resources are valuable, but access to clinical text is
constrained by privacy and institutional data-sharing restrictions
\cite{chapman2011}. Gold-standard construction can require detailed guidelines,
independent annotation, and adjudication \cite{uzuner2010annotation}. The rapid
adoption of LLMs creates an additional need for
benchmarks that are reproducible, extensible, and organized around interpretable
failure modes rather than a single aggregate test set.

Knowledge- and template-driven generators that do not ingest patient narratives
can broaden controlled benchmark coverage without directly distributing patient
records. Their value is greatest when the generation process is transparent,
task dimensions are controlled, and the output is carefully documented
\cite{walonoski2018,gebru2021}.
They do not replace real clinical corpora: template consistency does not by
itself establish clinical validity or reproduce the distribution of electronic
health records.

We present \dataset\ as a framework and a reference dataset instantiation. The
framework links normalized oncology knowledge inputs to a challenge taxonomy and
generates synthetic text together with structured targets. The present paper
makes four contributions:

\begin{enumerate}[leftmargin=*,nosep]
    \item a reproducible adapter from a pinned \kbadj{} relational snapshot to
    generator-facing drug--alias, regimen--component, and condition views,
    separated from author-specified configuration resources;
    \item a five-category, 20-subcategory taxonomy covering canonical
    mentions, prescribing attributes, regimen semantics, contextual status,
    and noisy surface forms;
    \item a generation engine that is deterministic under a fixed \kbadj{}
    revision, adapter, materialized views, code, software environment,
    configuration, ordering, and seed, with 347 templates, constrained entity
    sampling, and scenario-level text and annotation construction; and
    \item benchmark interfaces for LLM entity recognition, normalization,
    attribute extraction, contextual interpretation, and regimen resolution,
    with category-level analysis and extensible configurations.
\end{enumerate}

The framework reduces the amount of sample-by-sample labeling needed to create
an initial benchmark, but it does not remove the need for expert review. No
independent clinical validation was performed for the current release, and the
dataset is intended for extraction research rather than clinical decision
support, prescribing, or safety evaluation.

\section{Related Work and Knowledge Resources}
\label{sec:related}

\subsection{Medication Information Extraction}

The i2b2 medication challenge established a common evaluation setting for
extracting medication names and associated fields from discharge summaries
\cite{uzuner2010}. MADE 1.0 focused on medications, indications, and adverse
events in longitudinal EHR notes from patients with cancer
\cite{jagannatha2019}. The 2018 n2c2
shared task evaluated medication and adverse-event concepts and relations
\cite{henry2020}, whereas the 2022 n2c2 task emphasized contextualized
medication changes \cite{mahajan2023}. More recently, large language models
have been studied as zero- and few-shot clinical information extractors
\cite{agrawal2022}. Such systems make it useful to expose task-specific
benchmark slices that can distinguish surface normalization, attribute
recovery, contextual interpretation, and implicit regimen knowledge.

\subsection{Synthetic Clinical Data and Dataset Documentation}

Synthetic health data can facilitate method development and reproducible
testing when real records are inaccessible; Synthea is a prominent example of
a transparent generator for synthetic longitudinal records
\cite{walonoski2018}. Template-generated clinical text serves a different
purpose: it permits direct control over which linguistic and semantic challenge
appears in each sample and can emit structured labels at generation time.
Because synthetic data inherit the assumptions of their generators, clear
documentation of motivation, composition, construction, intended use, and
limitations remains essential \cite{gebru2021}. \dataset\ follows this
framework-oriented view: the generator is the primary contribution, and the
2,000-row collection is one deterministic output obtained when the input
content and ordering, code, configuration, and seed are fixed.

\subsection{Drug and Regimen Terminologies}

Normalized input views can be constructed from complementary resources such as
RxNorm, HemOnc, NCI Thesaurus, DrugBank, AACT, CanMED, and SEER*Rx
\cite{nelson2011,warner2019,sioutos2007,knox2024,tasneem2012,
canmed2024,seerrx2026}. For the reference release, these resources are
integrated through a pinned revision of \kbfull{}
publication build \cite{shivanna2026uotd}.
\ifanonymous
The identity-bearing build label and revision are withheld during review.
\else
The build is \code{publication-build-2026-08-09} at Git commit
\code{e4ba3722b550}.
\fi
\kbname{} applies deterministic
source-priority integration and a final RxNorm semantic gate, and publishes
relational drug, synonym, regimen, component, and condition tables. AACT is
retained by \kbname{} as a stand-alone clinical-trials table and is not consumed by
the \dataset\ adapter. \kbname{}'s published QA establishes file, schema,
referential, and identifier consistency; it does not clinically adjudicate
regimen component semantics.

\section{Framework Architecture}
\label{sec:framework}

The framework separates benchmark construction into four layers: knowledge,
configuration, generation, and delivery. This separation allows the input
vocabulary and benchmark taxonomy to evolve independently. A particular
collection such as \dataset\ is produced by fixing the \kbadj{} revision, adapter,
materialized-view content and ordering, code, software environment,
configuration, category targets, and random seed.

\subsection{Knowledge and Normalization Layer}

The released generator consumes three denormalized files under
\code{data/knowledge}. \code{drug\_table.csv} stores an upstream drug identifier,
canonical display name, and a JSON array of aliases;
\code{regimen\_table.csv} stores an upstream regimen identifier, name, projection
method, and a JSON array of component drugs; and
\code{Conditions\_And\_Regimens.csv} retains condition--regimen associations.
JSON arrays prevent commas inside names from being interpreted as list
delimiters. The loaders retain backward compatibility with the legacy
comma-separated representation, but the JSON-array views are the normative
release inputs.

The versioned adapter \code{scripts/build\_uotd\_inputs.py} verifies SHA-256
hashes for five production tables from the pinned \kbadj{} revision, normalizes
Unicode and whitespace, performs identifier-based joins, and orders rows and
list values deterministically. It writes complete input and output provenance
to \code{data/knowledge/provenance.json}. The views can be rebuilt from either
a local upstream checkout or the five production tables at the pinned public
commit.

Inspection of the upstream regimen relations identified source-integration cases
in which a direct union of linked components would over-expand a benchmark
target. The OncoRx adapter therefore applies a conservative, auditable
projection: it deduplicates direct links; retains component names only when
they occur as token-bounded terms in the canonical regimen name; requires at
least two retained components; excludes a row when an explicit
author-controlled drug surface in the regimen name is missing from the upstream
links; and adds a terminal parenthetical acronym only when it is release-unique
for one retained component tuple. \code{regimen\_projection\_audit.csv} records every retained
and quarantined upstream regimen and the reason for that decision. These rules are
benchmark-eligibility safeguards, not clinical validation of a regimen.

Upstream normalization reconciles source concepts, aliases, identifiers, and
provenance. The subsequent flattening retains only the fields required by the
generator. This avoids source-specific joins during generation and makes the
exact input content and ordering explicit for reproducibility.

Table~\ref{tab:source-roles} summarizes the role of each upstream knowledge
source. It reports functional roles rather than combining source row counts,
because releases, access terms, and filtering rules are source-specific.

\begin{table*}[t]
\caption{Roles of upstream resources in the pinned \kbadj{} build. The benchmark
adapter consumes the materialized production tables rather than querying these resources
during generation.}
\label{tab:source-roles}
\centering
\footnotesize
\begin{tabularx}{\textwidth}{@{}l Y Y@{}}
\toprule
\textbf{Resource} & \textbf{Primary framework role} & \textbf{Normalized contribution} \\
\midrule
RxNorm & Clinical drug normalization and synonym linkage & Canonical drug concepts, identifiers, and surface variants \\
HemOnc & Oncology regimen representation & Regimen names, aliases, and component agents \\
NCI Thesaurus & Cancer-oriented terminology & Disease, drug, therapy, and molecular terminology \\
DrugBank & Drug nomenclature & Alternative names and identifiers subject to source access terms \\
AACT & ClinicalTrials.gov intervention representation & Stand-alone upstream table; not consumed by the benchmark adapter \\
CanMED & Oncology medication coding & NDC nomenclature used in the pinned \kbadj{} build \\
SEER*Rx & Antineoplastic drug and regimen classification & Treatment categories, drug names, and regimen terminology \\
\bottomrule
\end{tabularx}
\end{table*}

\subsection{Configuration and Output Contracts}

The configuration layer defines the category hierarchy, target sample counts,
intended difficulty labels, random seed, curated oncology and supportive-care
pools, surface-variation maps, and attribute option lists. Templates use named
slots such as \code{drug1}, \code{regimen\_name}, \code{dose1},
\code{route1}, \code{freq1}, \code{condition}, and \code{intent}.
Scenario-specific functions encode sampled values selected as targets for that
subcategory; not every rendered slot is serialized as an annotation field.

The output contract is likewise source-independent. Every sample contains an
identifier, synthetic text, category metadata, a drug-object list, and a
regimen-object list. JSONL is the canonical nested representation; CSV export
serializes object lists as JSON strings for inspection and tool integration.
The JSONL/CSV delivery pattern can be retained across differently sized
benchmark versions and future domain extensions, although domain-specific
schemas and task manifests may change.

\begin{table}[t]
\caption{Layered contracts of the synthetic-benchmark framework.}
\label{tab:contracts}
\centering
\footnotesize
\begin{tabularx}{\columnwidth}{@{}l Y@{}}
\toprule
\textbf{Layer} & \textbf{Contract} \\
\midrule
Knowledge & Canonical drugs and aliases, regimen names and components, and condition terms \\
Configuration & Taxonomy, quotas, templates, candidate pools, seed, attribute options \\
Generation & Slot values, rendered text, task-selected drug/regimen objects \\
Delivery & JSONL/CSV schemas, category metadata, reusable benchmark subsets \\
\bottomrule
\end{tabularx}
\end{table}

\section{Reference Benchmark Construction}
\label{sec:design}

\subsection{Challenge Taxonomy}

Table~\ref{tab:taxonomy} summarizes the target distribution. The labels
\emph{Easy}, \emph{Medium}, \emph{Hard}, and \emph{Very Hard} are generator
designations, not human- or model-validated difficulty measurements. Categories
are diagnostic axes rather than mutually exclusive clinical phenomena.

\begin{table*}[t]
\caption{Taxonomy and target counts. Each category contains four
subcategories; parenthetical values are sample counts. Difficulty is an
intended generator label.}
\label{tab:taxonomy}
\centering
\footnotesize
\begin{tabularx}{\textwidth}{@{}c l Y r l@{}}
\toprule
\textbf{Cat.} & \textbf{Axis} & \textbf{Subcategories} & \textbf{$n$} & \textbf{Intended difficulty} \\
\midrule
C1 & Core extraction & Single canonical drug (150); drug with dose/route (200); two drugs (150); supportive care (100) & 600 & Easy--Medium \\
C2 & Prescribing information & Dose--route--frequency (150); taper/titration (100); prn/conditional (100); duration/stop (100) & 450 & Medium--Hard \\
C3 & Regimen semantics & Explicit multi-drug regimen (150); acronym only (150); partial component list (100); cycle/line/intent (50) & 450 & Hard--Very Hard \\
C4 & Context and safety language & Discontinued/held (100); allergy/ADR (80); negated (70); history/conflict (50) & 300 & Medium--Hard \\
C5 & Noise and ambiguity & Abbreviation (60); brand name (60); misspelling (50); distractor-rich paragraph (30) & 200 & Medium--Very Hard \\
\midrule
 & \textbf{Total} & 20 subcategories & \textbf{2,000} & \\
\bottomrule
\end{tabularx}
\end{table*}

The five categories are designed to isolate different construction axes. C1
varies the number and type of explicit medication mentions. C2 adds structured
administration language such as route, frequency, taper, conditional use, and
duration. C3 shifts the representation unit from individual drugs to named
regimens, including acronym-only and partial-component descriptions. C4 places
otherwise recognizable medications under discontinuation, history, negation,
allergy, or conflict language. C5 changes the observed surface through curated
abbreviations, brands, misspellings, and longer distractor text. These axes are
not mutually exclusive clinical phenomena; they are template-generation
controls used to organize the artifact.

Each subcategory target exceeds the size of its template pool, so templates are
sampled repeatedly during generation. Table~\ref{tab:template-allocation}
reports the target divided by the number of available templates. The ratio
ranges from 3.3 to 10.0 but does not imply equal use of individual templates.

\begin{table*}[t]
\caption{Template-pool sizes in the supplied \code{templates.py}. Ratio is the
configured target divided by the number of available templates; random
sampling does not enforce equal use of individual templates.}
\label{tab:template-allocation}
\centering
\footnotesize
\begin{tabular}{@{}c r r r @{\hspace{1.5em}} c r r r@{}}
\toprule
\textbf{Subcat.} & \textbf{Target} & \textbf{Templates} & \textbf{Ratio} &
\textbf{Subcat.} & \textbf{Target} & \textbf{Templates} & \textbf{Ratio} \\
\midrule
C1.1 & 150 & 30 & 5.0 & C3.3 & 100 & 15 & 6.7 \\
C1.2 & 200 & 25 & 8.0 & C3.4 &  50 & 10 & 5.0 \\
C1.3 & 150 & 25 & 6.0 & C4.1 & 100 & 20 & 5.0 \\
C1.4 & 100 & 20 & 5.0 & C4.2 &  80 & 16 & 5.0 \\
C2.1 & 150 & 20 & 7.5 & C4.3 &  70 & 20 & 3.5 \\
C2.2 & 100 & 15 & 6.7 & C4.4 &  50 & 13 & 3.8 \\
C2.3 & 100 & 15 & 6.7 & C5.1 &  60 & 15 & 4.0 \\
C2.4 & 100 & 15 & 6.7 & C5.2 &  60 & 15 & 4.0 \\
C3.1 & 150 & 15 & 10.0 & C5.3 &  50 & 15 & 3.3 \\
C3.2 & 150 & 20 & 7.5 & C5.4 &  30 &  8 & 3.8 \\
\midrule
\multicolumn{2}{@{}l}{\textbf{Total}} & \textbf{347} & &
\multicolumn{2}{l}{\textbf{Configured samples}} & \textbf{2,000} & \\
\bottomrule
\end{tabular}
\end{table*}

\subsection{Generation Inputs}

The release requires all three checked-in views: canonical drugs with aliases,
benchmark-eligible regimens with projected component drugs, and
condition--regimen associations. Missing, malformed, or checksum-inconsistent
inputs stop the release workflow rather than silently selecting a fallback
vocabulary. The final materialization contains
\FinalDrugViewRows{} drug rows,
\FinalRegimenViewRows{} regimen rows, and
\FinalConditionViewRows{} condition--regimen rows; the complete
counts and hashes are recorded in the machine-readable provenance manifest.

The supplied reference configuration additionally contains author-specified
control resources: 122 oncology-oriented drug names, 42 supportive-care drugs,
43 per-drug dose-option profiles, 57 brand-to-generic mappings, 37
abbreviation mappings, 20 misspelling dictionaries with four variants per
drug, 20 drug-specific adverse-reaction profiles, 15 prn condition profiles,
and four anthracycline cumulative-dose scenario constraints. The candidate
pools keep the generator focused on oncology and supportive-care language while
the input views retain a broader normalized vocabulary. The dose, route,
frequency, taper, adverse-reaction, cumulative-dose, and drug--reason values
are generator constraints, not independently validated clinical
recommendations. They can be replaced with versioned, clinically reviewed
joint tuples without changing the template or serialization layers.

Table~\ref{tab:controls} summarizes the controls in the reference
configuration. They are explicit configuration objects rather than values
embedded inside prompt text, so they can be versioned, expanded, or replaced
independently of the template library.

\begin{table*}[t]
\caption{Author-specified control resources in the reference configuration.
They constrain generation and are not presented as clinically validated
recommendations.}
\label{tab:controls}
\centering
\footnotesize
\begin{tabularx}{\textwidth}{@{}l r Y@{}}
\toprule
\textbf{Control resource} & \textbf{Size} & \textbf{Framework function} \\
\midrule
Oncology-oriented pool & 122 & Candidate sampling across treatment and supportive contexts \\
Supportive-care pool & 42 & Antiemetic, growth-factor, analgesic, and related scenarios \\
Drug-specific dose profiles & 43 & Candidate dose, route, and frequency values \\
Brand-to-generic mappings & 57 & Controlled brand normalization challenges \\
Abbreviation mappings & 37 & Controlled abbreviation normalization challenges \\
Misspelling dictionaries & 20 drugs $\times$ 4 & Seeded, controlled surface corruption \\
Adverse-reaction profiles & 20 drugs & Allergy and adverse-reaction context templates \\
prn condition profiles & 15 drugs & Drug--reason pairing for a profiled subset of conditional-use templates \\
Cumulative-dose constraints & 4 drugs & Anthracycline-specific values in cumulative-dose scenarios \\
\bottomrule
\end{tabularx}
\end{table*}

The reference configuration uses ten note-type labels. Each subcategory has an
author-specified eligible subset selected for compatibility with its template
family, and selection within that subset uses the same seeded pseudorandom
generator as the remaining scenario values. These labels do not estimate the
prevalence of document types in an external EHR corpus and are not independently
validated note classifications.

The authors report that Anthropic Claude Sonnet 4.7 assisted initial drafting
of portions of the source code, template library, configuration resources, and
explanatory text \cite{anthropic2026}. The exact prompts and decoding settings
were not retained, so that developmental drafting process cannot be replayed
exactly. The retained, versioned artifacts were subsequently audited and
revised. No LLM is invoked while rendering individual dataset rows; release
generation uses the deterministic Python procedure described below.

\subsection{Generation Procedure}

Figure~\ref{fig:pipeline} shows the benchmark-construction path. With random
seed 42, the reference configuration iterates through 20 subcategory
specifications, selects a template from the corresponding pool, samples
entities and optional attributes, renders text, and constructs structured
annotations. Duplicate detection uses a lowercase, whitespace-trimmed key while
preserving the rendered text and permits up to 100 candidate-generation
attempts. For regimen-bearing scenarios that require a condition, the
condition is sampled from associations for the selected upstream regimen rather
than from an independent global list. The exporter serializes object lists into CSV columns
and constructs an exact 80/20 split within every subcategory while keeping each
source-template group wholly in train or test.

\begin{figure*}[t]
\centering
\begin{tikzpicture}[
  stage/.style={draw=black!65, rounded corners=2pt, align=center,
    text width=2.85cm, minimum height=1.55cm, inner sep=3pt,
    font=\scriptsize, fill=blue!4},
  arrow/.style={-{Stealth[length=5pt]}, semithick, draw=black!70},
  node distance=0.43cm]
\node[stage] (inputs) {\textbf{1. Knowledge materialization}\\Pinned upstream tables\\SHA-verified adapter\\Three generator views};
\node[stage, right=of inputs] (constraints) {\textbf{2. Curated lists}\\Oncology/supportive pools\\Alias and typo maps\\Dose option profiles};
\node[stage, right=of constraints] (instantiate) {\textbf{3. Instantiation}\\20 subcategories\\347 templates\\Seeded sampling};
\node[stage, right=of instantiate] (labels) {\textbf{4. Scenario output}\\Synthetic text\\Drug/regimen objects\\Status/prescribing fields};
\node[stage, right=of labels] (deliver) {\textbf{5. Audit and release}\\Nested-schema checks\\Grounding and offsets\\JSONL/CSV and hashes};
\draw[arrow] (inputs) -- (constraints);
\draw[arrow] (constraints) -- (instantiate);
\draw[arrow] (instantiate) -- (labels);
\draw[arrow] (labels) -- (deliver);
\end{tikzpicture}
\caption{Knowledge-grounded synthetic benchmark generation pipeline. A fixed
upstream revision, adapter, normalized-view content and ordering, author-specified
constraints, template targets, software environment, code, and seed produce a
deterministic dataset output.}
\label{fig:pipeline}
\end{figure*}

Generation is target-driven rather than distribution-estimating. For each
subcategory, the implementation constructs the configured number of valid
samples or stops with an error. Identifiers use a monotonically increasing
counter. Text and task-specific target objects are constructed from shared
sampled slot values; exceptions do not produce simplified rows with targets
from a failed scenario.

The duplicate-control loop forms a lowercased, trimmed comparison key from the
generated text and permits up to 100 candidate-generation attempts. If every
attempt collides, generation stops rather than admitting a known duplicate.
Runs are repeatable when the pinned upstream revision, adapter, materialized-view
content and ordering, code, software environment, configuration, and seed are
fixed.

The companion integrity script checks required and optional field types,
nested objects, enumerations, category quotas, difficulty and note-type values,
unique sample identifiers, generated-drug counts, and character-offset/surface
agreement. It verifies normalized drugs, regimen names, and projected regimen
components against the canonical generator views; flags unresolved template
slots and normalized duplicate text; and exits nonzero on failure. Split checks
confirm exhaustive, disjoint identifiers and zero shared source templates.
These controls establish structural,
referential, and release consistency only; they do not establish joint clinical
correctness or exhaustive annotation of incidental medication language. The
exporter serializes nested objects as JSON strings and creates human-readable
summaries.

\paragraph{Observed integrity results.}
For the final 2,000-row release, the release checks report
\FinalValidationErrors{} structural or referential errors and
\FinalDuplicateTexts{} normalized duplicate texts from the strict validator;
the split manifest and exporter checks report
\FinalSplitLeakageErrors{} source-template split-overlap errors. No automated
result in this paragraph constitutes clinical validation.

\subsection{Reference Implementation}

The implementation mirrors the layered design rather than placing generation
logic in one monolithic script. Table~\ref{tab:modules} identifies the core
modules and release utilities. This decomposition makes it possible to
modify a taxonomy, add templates, revise schema fields, or change export policy
without rewriting the other layers. \ifanonymous
Source code, validation utilities, and release materials are described by an
artifact citation withheld from the paper body for review
\cite{shivanna2026github}.\else
Source code, validation utilities, and release materials are available in the
project repository \cite{shivanna2026github}.\fi

\begin{table*}[t]
\caption{Reference software modules and framework responsibilities.}
\label{tab:modules}
\centering
\footnotesize
\begin{tabularx}{\textwidth}{@{}l Y@{}}
\toprule
\textbf{Module} & \textbf{Responsibility} \\
\midrule
\code{scripts/build\_uotd\_inputs.py} & Pinned upstream acquisition, checksum verification, deterministic view materialization, regimen projection, and provenance/audit output \\
\code{dataset\_config.py} & Seed, category quotas, candidate vocabularies, attribute profiles, and controlled variation maps \\
\code{templates.py} & 347 template strings organized by 20 subcategories \\
\code{generate\_dataset.py} & Knowledge-view loaders, samplers, subcategory generators, text--target construction, and JSONL serialization \\
\code{schema.py} & Drug, prescribing-information, regimen, offset, and benchmark-sample output objects \\
\code{validate\_dataset.py} & Nested schema, enum, grounding, offset, evidence, quota, and duplicate checks with nonzero failure status \\
\code{export\_csv.py} & JSONL/CSV export and deterministic subcategory-stratified, source-template-grouped partitioning \\
\code{scripts/reproduce\_release.py} & One-command generation, validation, export, profiling, manifest construction, and byte comparison \\
\code{scripts/profile\_release.py} & Dataset statistics and generated LaTeX values used by the paper tables \\
\code{data/knowledge/provenance.json} & Upstream revision, source and output hashes, view counts, transformation policy, and projection summary \\
\code{data/knowledge/regimen\_projection\_audit.csv} & Per-regimen retained/quarantined status and component-projection evidence \\
\bottomrule
\end{tabularx}
\end{table*}

\paragraph{Reproducible release.}
The command \code{python3 scripts/reproduce\_release.py {-}{-}check} verifies
the pinned knowledge views, regenerates JSONL in a temporary directory, runs
the strict validator, rebuilds the full and split JSONL/CSV artifacts and paper
statistics, and byte-compares every result with the archive. The reference
artifact uses CPython 3.9.6 with no external Python packages. Its canonical
JSONL SHA-256 is
\code{\FinalJsonlShaA}\allowbreak\code{\FinalJsonlShaB}. The release manifest
records complete source and artifact hashes. For double-anonymous review,
identity-bearing repository and dataset revisions are withheld from the paper
body and should be restored in the camera-ready version.

\subsection{Annotation Objects}

Each row contains an identifier, synthetic text, category, subcategory,
intended difficulty, auxiliary note-type label, generated-drug count,
drug-object list, and regimen-object list.

A drug object contains a stored surface string, normalized name, inclusive
start and exclusive end character offsets, medication state, separate
contextual flags, an optional reason or adverse event, and optional prescribing
information. The field \code{evidence\_type} distinguishes a literal
\code{explicit\_surface} from a \code{regimen\_inference} whose evidence span
is the named regimen. The field \code{status} records temporal/action state and takes
one of \code{current}, \code{planned}, \code{historical},
\code{discontinued}, \code{hold}, or \code{unknown}. The Boolean fields
\code{negated}, \code{allergy}, and \code{uncertain} are separate contextual
properties rather than additional status values. The separate
\code{adverse\_event} field stores a rendered reaction or toxicity when the
scenario attributes it to the medication.

The serialized field \code{sig} stores optional prescribing information,
including dose value and unit, route, frequency, duration, form, \code{prn},
taper, cycle-day, and infusion-time fields. The reference release populates
only fields supported by the rendered text. A regimen object contains its
surface and normalized forms, inclusive/exclusive character offsets,
benchmark-eligible normalized components, and optional cycle and intent
metadata. The optional \code{intent} field is separate from medication status
and may take one of \code{neoadjuvant}, \code{adjuvant}, \code{first\_line},
\code{second\_line}, \code{third\_line\_plus}, \code{maintenance},
\code{palliative}, \code{curative}, \code{salvage}, \code{conditioning},
\code{consolidation}, \code{induction}, or \code{other}. Cycle and intent
metadata are serialized only when the corresponding value is expressed in the
synthetic text.

The schema deliberately supports both literal medication mentions and regimen
resolution. A named regimen can therefore carry a normalized component list
even when the individual agents are not written in the snippet. Category and
subcategory metadata can be combined with a task manifest to identify eligible
rows and select the target fields relevant to an LLM benchmarking task. For a
component inferred from a named regimen, any stored surface and offsets identify
the regimen evidence span; they do not assert a literal span for the unwritten
component name.

\section{Generated Dataset and Benchmark Interfaces}
\label{sec:dataset}

\subsection{Reference Dataset Composition}

Binding the framework to the taxonomy, template library, curated lists, and
seed described above produces the 2,000-row reference dataset summarized in
Table~\ref{tab:profile}. These statistics characterize the generated resource
rather than a model experiment. The collection contains
\FinalDrugObjects{} structured drug objects and
\FinalRegimenObjects{} regimen objects, enabling one dataset to
support multiple extraction interfaces. \ifanonymous
The corresponding data artifact is identified through an anonymous artifact
citation \cite{shivanna2026hf}.\else
The corresponding dataset release is available through Hugging Face
\cite{shivanna2026hf}.\fi

\begin{table}[t]
\caption{Composition of the reference \dataset\ configuration.}
\label{tab:profile}
\centering
\footnotesize
\begin{tabularx}{\columnwidth}{@{}Y r@{}}
\toprule
\textbf{Statistic} & \textbf{Value} \\
\midrule
Samples / configured note types & 2,000 / 10 \\
Categories / subcategories & 5 / 20 \\
Templates & 347 \\
Drug annotation objects & \FinalDrugObjects \\
Unique normalized drugs / stored evidence surfaces & \FinalUniqueDrugs{} / \FinalUniqueEvidenceSurfaces \\
Literal drug / regimen-evidence surfaces & \FinalUniqueExplicitSurfaces{} / \FinalUniqueInferenceSurfaces \\
Drug objects with prescribing information & \FinalSigObjects \\
Regimen objects / unique names & \FinalRegimenObjects{} / \FinalUniqueRegimens \\
Component assignments / unique components & \FinalComponentAssignments{} / \FinalUniqueComponents \\
Mean / median whitespace terms & \FinalMeanTerms{} / \FinalMedianTerms \\
Easy / Medium / Hard / Very Hard & 350 / 840 / 580 / 230 \\
\bottomrule
\end{tabularx}
\end{table}

Table~\ref{tab:category-profile} shows how annotation structures follow the
five framework axes. C1 and C2 contribute prescribing-information objects, C3
concentrates regimen structures, C4 concentrates contextual medication states,
and C5 applies surface perturbations and distractors. The allocation is
intentional: benchmark users
can select a category aligned with the extraction capability being studied
instead of treating every row as the same task.

\begin{table*}[t]
\caption{Reference-dataset profile by framework category. Mean terms are
computed by whitespace splitting.}
\label{tab:category-profile}
\centering
\footnotesize
\begin{tabular}{@{}c r r r r r r@{}}
\toprule
\textbf{Category} & \textbf{Rows} & \textbf{Template pool} &
\textbf{Mean terms} & \textbf{Drug objects} & \textbf{Prescr. objects} &
\textbf{Regimen objects} \\
\midrule
C1 & 600 & 100 & \FinalCOneMeanTerms & \FinalCOneDrugs & \FinalCOneSigs & \FinalCOneRegimens \\
C2 & 450 & 65  & \FinalCTwoMeanTerms & \FinalCTwoDrugs & \FinalCTwoSigs & \FinalCTwoRegimens \\
C3 & 450 & 60  & \FinalCThreeMeanTerms & \FinalCThreeDrugs & \FinalCThreeSigs & \FinalCThreeRegimens \\
C4 & 300 & 69  & \FinalCFourMeanTerms & \FinalCFourDrugs & \FinalCFourSigs & \FinalCFourRegimens \\
C5 & 200 & 53  & \FinalCFiveMeanTerms & \FinalCFiveDrugs & \FinalCFiveSigs & \FinalCFiveRegimens \\
\midrule
\textbf{Total} & \textbf{2,000} & \textbf{347} & \FinalMeanTerms &
\FinalDrugObjects & \FinalSigObjects & \FinalRegimenObjects \\
\bottomrule
\end{tabular}
\end{table*}

\subsection{Illustrative Generated Instances}

Table~\ref{tab:examples} illustrates the coupling between rendered text and
structured targets. The examples span two dose--route--frequency orders, a
partial regimen description, a negated medication,
and a misspelled surface form. They are examples of framework behavior and
have not undergone independent clinical review. The rows were selected from
and checked against the final regenerated JSONL artifact.

\begin{table*}[t]
\caption{Representative synthetic rows and generated target structures.}
\label{tab:examples}
\centering
\footnotesize
\begin{tabularx}{\textwidth}{@{}l l Y Y@{}}
\toprule
\textbf{ID} & \textbf{Subcat.} & \textbf{Synthetic text} &
\textbf{Generated target summary} \\
\midrule
ONCORX-0603 & C2.1 & Current regimen: Bevacizumab 15 mg/kg IV q3w. & Bevacizumab; dose 15 mg/kg; route IV; frequency q3w \\
ONCORX-0606 & C2.1 & Current regimen: Capecitabine 850 mg/m2 PO BID days 1-14 q3w. & Capecitabine; dose 850 mg/m2; route PO; frequency BID days 1-14 q3w \\
ONCORX-1354 & C3.3 & cisplatin and etoposide (ep) and bevacizumab: Etoposide IV, with remaining agents scheduled separately. & Regimen components: Bevacizumab, Cisplatin, Etoposide; Etoposide route IV \\
ONCORX-1681 & C4.3 & No indication for Tamoxifen at this time. & Tamoxifen; negated=true; status=unknown \\
ONCORX-1922 & C5.3 & Infused paclitaxle 175 mg/m2 IV. & paclitaxle $\rightarrow$ Paclitaxel; dose 175 mg/m2; route IV \\
\bottomrule
\end{tabularx}
\end{table*}

\subsection{LLM Benchmarking Interface}

Prior work has evaluated zero- and few-shot LLMs for clinical information
extraction and implemented structured LLM extraction pipelines
\cite{agrawal2022,hsu2025}. The schema developed here is intended to support
future zero-shot, few-shot, retrieval-augmented, and fine-tuned evaluations. No
LLM results are reported in this paper.

Table~\ref{tab:llm-tasks} maps eligible framework categories to four extraction
interfaces. Separating the interfaces avoids requiring a model to emit fields
that a category was not constructed to represent. It also supports
category-level reporting: for example, normalization behavior can be examined
separately for canonical names, brands, abbreviations, and misspellings, while
regimen resolution can be isolated from literal drug recognition. A future
evaluation release should add a versioned task manifest identifying eligible
rows and the target projection for each interface.

\begin{table*}[t]
\caption{Intended LLM benchmarking interfaces supported by the framework.}
\label{tab:llm-tasks}
\centering
\footnotesize
\begin{tabularx}{\textwidth}{@{}l l Y Y@{}}
\toprule
\textbf{Interface} & \textbf{Primary categories} &
\textbf{Requested model output} & \textbf{Controlled challenge} \\
\midrule
Mention extraction and normalization & C1, C5.1--C5.3 &
Drug surface and canonical generic name &
Canonical names, brands, abbreviations, and misspellings \\
Prescribing-information extraction & C1.2; eligible C2 rows &
Drug-linked dose, unit, route, frequency, duration, \code{prn}, and taper fields when populated &
Attribute combinations and order phrasing \\
Regimen resolution & C3.2--C3.4; eligible C3.1 rows &
Regimen name and benchmark-eligible component projection, with cycle and intent fields when visibly supported &
Explicit, acronym-only, and partial-component descriptions \\
Contextual medication interpretation & C4 &
Medication temporal/action status plus separate negation and allergy flags &
Current, planned, historical, held, discontinued, negated, and allergy contexts \\
\bottomrule
\end{tabularx}
\end{table*}

A reproducible LLM study built on the framework should freeze the upstream
revision, adapter and knowledge views, configuration, template version, seed,
and task schema. It should also
state whether the model can access an external drug or regimen knowledge base,
whether decoding is constrained to valid JSON, and whether demonstrations are
selected from the same category. Future template-, entity-, and
regimen-disjoint configurations can support stronger generalization studies.
These are benchmark-design choices, not properties that must remain fixed
across all uses of the framework. Future paired configurations can hold the
target entity fixed while varying surface form, syntax, context, or noise for
controlled robustness analysis.

\subsection{Extensibility and Scale}

The framework is modular along three dimensions. First, a domain extension can
reuse orchestration and delivery while supplying domain-specific knowledge
views, schema objects, templates, candidate pools, and generator functions.
Second, new task axes can be registered without changing existing categories.
Third, targets and seeds can produce benchmark versions with different sizes
and challenge mixtures.

The reference implementation uses one pseudorandom generator and constructs
categories sequentially. Parallel generation is an architectural extension,
not a property of the current code. A future distributed implementation could
assign deterministic seed ranges to shards and reconcile duplicates globally;
reproducibility and scaling would need to be measured before performance claims
are made on high-performance computing systems.

\section{Discussion}
\label{sec:discussion}

\subsection{Why a Generation Framework is Needed}

Access to clinical text remains limited by privacy, institutional sharing
restrictions, and the effort required to create expert ground truth
\cite{chapman2011,uzuner2010annotation}. Medication and regimen extraction add
specialized labels that often require clinical knowledge, making it difficult
to build a separate human-annotated corpus for every task variant. Synthetic
corpora have therefore been investigated as complementary resources for
clinical entity extraction and, in some settings, relation extraction
\cite{brekke2021,li2021}.

The framework addresses this bottleneck by producing text and task-selected
targets within the same scenario-specific generation function from shared
sampled values. This reduces the initial sample-by-sample labeling burden and
provides configured control over category
targets and surface perturbations; the reference instantiation realizes all
configured category counts. An investigator can add a failure mode, regenerate
a slice from fixed inputs, and preserve a common output schema. The benefit is
controlled coverage, not a claim of real-note representativeness.

\subsection{Knowledge Grounding and Controlled Variation}

Knowledge grounding separates what an entity means from how it appears in text.
Canonical drug names, projected regimen components, and regimen--condition
associations enter through the pinned knowledge-derived views, whereas brands,
abbreviations, misspellings, and selected drug--reason relations enter through
explicit author-controlled configuration maps. Although the loader constructs
a general alias map, current generators use controlled surface-variation maps
rather than sampling the entire alias inventory. Templates then vary syntax
and context while retaining task-specific targets. Referential grounding and
lexical projection improve traceability, but neither establishes clinical
correctness. Recent work on knowledge-guided clinical text generation supports
this general design direction \cite{xu2024}.

This separation is useful for LLM evaluation. Direct models can be compared
with models that receive a drug or regimen knowledge source, and constrained
JSON generation can be compared with free-form extraction under the same
category definitions. Fine-grained categories help identify whether a failure
arises from mention recognition, normalization, contextual interpretation, or
knowledge-dependent regimen expansion rather than collapsing all behaviors into
one score.

\subsection{Complementarity with Human-Labeled Data}

Synthetic benchmarks should be used alongside, not instead of, expert-labeled
clinical corpora. Studies of synthetic clinical notes have shown that they can
support entity-extraction development, while also emphasizing the importance of
testing transfer to real text \cite{brekke2021,li2021}. For \dataset, a sensible
future program is to use the synthetic categories for controlled stress testing
and use access-controlled real-note datasets for external validity. This
division preserves the control and configurability of the framework without
overstating the clinical realism of its outputs.
\section{Limitations and Future Work}
\label{sec:limitations}

\textbf{Synthetic scope.}
The reference dataset contains short, English-language, template-generated
snippets. It does not reproduce document-level coreference, longitudinal
chronology, institutional writing styles, copy-forward behavior, or the full
distribution of errors in real electronic health records. The intended
difficulty tiers are design labels rather than empirically established measures
of human or model difficulty. The scenario-conditioned note-type labels are
generator metadata and should not be interpreted as an estimate of document
prevalence or as validated note classification.

\textbf{Knowledge grounding and clinical validity.}
The framework obtains normalization targets from input tables and curated
mappings, but entity-level grounding does not establish the
clinical appropriateness of every disease, drug, dose, route, frequency, and
schedule combination. No independent clinical validation was performed for the
current dataset. It must not be used for prescribing, regimen validation,
contraindication assessment, or patient-care decisions. Future configurations
should use clinically reviewed joint tuples and a documented expert sampling
and adjudication protocol. The conservative OncoRx regimen projection prevents
known source-link unions from being copied directly into benchmark targets, but
it is a lexical, fail-closed eligibility rule. It can exclude valid regimens and
does not prove that every retained component set is clinically complete or
appropriate.

\textbf{Annotation and partition design.}
The current release provides sample-level, task-selective structured targets;
it does not guarantee exhaustive annotation of every medication-like string in
every category. Distractor-rich C5.4 templates may contain incidental medication
text outside the task-selected target list. The release provides character
offsets for stored evidence surfaces and an \code{evidence\_type} field that
separates literal drug surfaces from components inferred through a regimen
span, but it does not yet provide explicit task-eligibility masks or per-object
source-provenance identifiers. The supplied 80/20 split is an in-distribution,
subcategory-stratified, source-template-disjoint split; entity- and
regimen-disjoint partitions should complement it in future generalization
studies.

\textbf{Reproducibility and licensing.}
The reference release freezes the upstream revision and source-table hashes,
adapter, materialized views and their hashes, code, configuration, software
environment, random seed, generated artifacts, and validation report. The
benchmark can therefore be regenerated from the checked-in views, and the views
can be independently rebuilt from the pinned public production tables.
Reconstructing \kbname{} itself from raw third-party sources still requires users to
obtain license-appropriate snapshots; public availability of a derived file
does not supersede upstream access or redistribution terms. The upstream source
manifest, OncoRx provenance manifest, and regimen-projection audit document this
boundary. The generator does not require patient narratives, but governance of
source data and generated releases remains necessary.

\textbf{Future validation and benchmarking.}
The release audit covers nested schema and enumeration conformance,
surface--offset agreement, referential grounding, configured distributions,
duplicate control, partition exhaustiveness, and source-template overlap.
These automated checks do not establish clinical semantic consistency,
exhaustive task annotation, or transfer to real notes. Future work will add
documented expert review of a representative sample. A separate empirical study
will compare direct, few-shot, fine-tuned, and knowledge-augmented LLM
extractors using frozen prompts and model versions. Human evaluation of clinical
LLM outputs has included factuality, reasoning, bias, and the likelihood and
severity of potential harm \cite{singhal2023}.

\section{Conclusion}
\label{sec:conclusion}

\dataset\ presents synthetic clinical benchmark construction as a reusable
framework rather than a one-time corpus-generation exercise. Modular knowledge
inputs, a 20-subcategory task taxonomy, category-specific samplers, controlled
variation resources, 347 templates, and scenario-level text and target
generation support deterministic construction of structured benchmarks when
the upstream revision, adapter, materialized inputs, software environment, code,
configuration, ordering, and seed are fixed. The reference configuration
produces 2,000 examples spanning canonical entities, prescribing attributes, regimen resolution,
contextual mentions, and perturbed surface forms.

Such controlled synthetic datasets can complement scarce expert-annotated
clinical corpora by providing repeatable challenge sets for LLM-based entity
extraction and category-specific analysis. They do not replace evaluation on
real clinical text or establish clinical correctness. Future work will add
independent clinical review, exhaustive task annotations, broader language
generation, stronger partitioning, and systematic LLM comparisons.

\section*{Acknowledgment}

Anthropic Claude Sonnet 4.7 assisted initial drafts of portions of the code,
templates, configuration, and explanatory text, as disclosed in
Section~\ref{sec:design} \cite{anthropic2026}. OpenAI Codex was used to audit
and materially revise the knowledge adapter, generator, schema, validator,
splitter, tests, release artifacts, documentation, and manuscript
\cite{openai2026codex}. AI-assisted wording may appear in the abstract and
Sections~I--VIII. No LLM generated individual dataset rows at release time.
The authors reviewed and approved all retained code, data, and manuscript
revisions and remain fully responsible for the work's content. \ifanonymous
Author-identifying funding and contract acknowledgments
are omitted for anonymous review.\else
\blfootnote{Notice: This manuscript has been authored by UT-Battelle LLC under
contract DE-AC05-00OR22725 with the US Department of Energy (DOE). The US
government retains and the publisher, by accepting the article for publication,
acknowledges that the US government retains a nonexclusive, paid-up,
irrevocable, worldwide license to publish or reproduce the published form of
this manuscript, or allow others to do so, for US government purposes. DOE will
provide public access to these results of federally sponsored research in
accordance with the DOE Public Access Plan
(\url{https://www.energy.gov/doe-public-access-plan}).}
\fi

% Apply \IEEEtriggeratref only after the final text and references are frozen.
\begin{thebibliography}{99}

\bibitem{uzuner2010}
O. Uzuner, I. Solti, and E. Cadag, ``Extracting medication information from
clinical text,'' \emph{J. Am. Med. Inform. Assoc.}, vol. 17, no. 5,
pp. 514--518, 2010, doi: \href{https://doi.org/10.1136/jamia.2010.003947}{10.1136/jamia.2010.003947}.

\bibitem{jagannatha2019}
A. Jagannatha, F. Liu, W. Liu, and H. Yu, ``Overview of the first natural
language processing challenge for extracting medication, indication, and
adverse drug events from electronic health record notes (MADE 1.0),''
\emph{Drug Saf.}, vol. 42, no. 1, pp. 99--111, 2019, doi:
\href{https://doi.org/10.1007/s40264-018-0762-z}{10.1007/s40264-018-0762-z}.

\bibitem{mahajan2023}
D. Mahajan, J. J. Liang, C.-H. Tsou, and O. Uzuner, ``Overview of the 2022
n2c2 shared task on contextualized medication event extraction in clinical
notes,'' \emph{J. Biomed. Inform.}, vol. 144, Art. no. 104432, 2023, doi:
\href{https://doi.org/10.1016/j.jbi.2023.104432}{10.1016/j.jbi.2023.104432}.

\bibitem{chapman2011}
W. W. Chapman, P. M. Nadkarni, L. Hirschman, L. W. D'Avolio, G. K. Savova,
and O. Uzuner, ``Overcoming barriers to NLP for clinical text: The role of
shared tasks and the need for additional creative solutions,'' \emph{J. Am.
Med. Inform. Assoc.}, vol. 18, no. 5, pp. 540--543, 2011, doi:
\href{https://doi.org/10.1136/amiajnl-2011-000465}{10.1136/amiajnl-2011-000465}.

\bibitem{uzuner2010annotation}
O. Uzuner, I. Solti, F. Xia, and E. Cadag, ``Community annotation experiment
for ground truth generation for the i2b2 medication challenge,'' \emph{J. Am.
Med. Inform. Assoc.}, vol. 17, no. 5, pp. 519--523, 2010, doi:
\href{https://doi.org/10.1136/jamia.2010.004200}{10.1136/jamia.2010.004200}.

\bibitem{walonoski2018}
J. Walonoski \emph{et al.}, ``Synthea: An approach, method, and software
mechanism for generating synthetic patients and the synthetic electronic
health care record,'' \emph{J. Am. Med. Inform. Assoc.}, vol. 25, no. 3,
pp. 230--238, 2018, doi:
\href{https://doi.org/10.1093/jamia/ocx079}{10.1093/jamia/ocx079}.

\bibitem{gebru2021}
T. Gebru \emph{et al.}, ``Datasheets for datasets,'' \emph{Commun. ACM},
vol. 64, no. 12, pp. 86--92, 2021, doi:
\href{https://doi.org/10.1145/3458723}{10.1145/3458723}.

\bibitem{henry2020}
S. Henry, K. Buchan, M. Filannino, A. Stubbs, and O. Uzuner, ``2018 n2c2
shared task on adverse drug events and medication extraction in electronic
health records,'' \emph{J. Am. Med. Inform. Assoc.}, vol. 27, no. 1,
pp. 3--12, 2020, doi:
\href{https://doi.org/10.1093/jamia/ocz166}{10.1093/jamia/ocz166}.

\bibitem{agrawal2022}
M. Agrawal, S. Hegselmann, H. Lang, Y. Kim, and D. Sontag, ``Large language
models are few-shot clinical information extractors,'' in \emph{Proc. 2022
Conf. Empirical Methods Natural Language Processing}, 2022, pp. 1998--2022,
doi: \href{https://doi.org/10.18653/v1/2022.emnlp-main.130}{10.18653/v1/2022.emnlp-main.130}.

\bibitem{nelson2011}
S. J. Nelson, K. Zeng, J. Kilbourne, T. Powell, and R. Moore, ``Normalized
names for clinical drugs: RxNorm at 6 years,'' \emph{J. Am. Med. Inform.
Assoc.}, vol. 18, no. 4, pp. 441--448, 2011, doi:
\href{https://doi.org/10.1136/amiajnl-2011-000116}{10.1136/amiajnl-2011-000116}.

\bibitem{warner2019}
J. L. Warner \emph{et al.}, ``HemOnc: A new standard vocabulary for
chemotherapy regimen representation in the OMOP common data model,''
\emph{J. Biomed. Inform.}, vol. 96, Art. no. 103239, 2019, doi:
\href{https://doi.org/10.1016/j.jbi.2019.103239}{10.1016/j.jbi.2019.103239}.

\bibitem{sioutos2007}
N. Sioutos, S. de Coronado, M. W. Haber, F. W. Hartel, W.-L. Shaiu, and
L. W. Wright, ``NCI Thesaurus: A semantic model integrating cancer-related
clinical and molecular information,'' \emph{J. Biomed. Inform.}, vol. 40,
no. 1, pp. 30--43, 2007, doi:
\href{https://doi.org/10.1016/j.jbi.2006.02.013}{10.1016/j.jbi.2006.02.013}.

\bibitem{knox2024}
C. Knox \emph{et al.}, ``DrugBank 6.0: The DrugBank knowledgebase for 2024,''
\emph{Nucleic Acids Res.}, vol. 52, no. D1, pp. D1265--D1275, 2024, doi:
\href{https://doi.org/10.1093/nar/gkad976}{10.1093/nar/gkad976}.

\bibitem{tasneem2012}
A. Tasneem \emph{et al.}, ``The database for aggregate analysis of
ClinicalTrials.gov (AACT) and subsequent regrouping by clinical specialty,''
\emph{PLOS ONE}, vol. 7, no. 3, Art. no. e33677, 2012, doi:
\href{https://doi.org/10.1371/journal.pone.0033677}{10.1371/journal.pone.0033677}.

\bibitem{canmed2024}
Division of Cancer Control and Population Sciences, National Cancer Institute,
``Cancer Medications Enquiry Database (CanMED),'' Surveillance Research
Program SEER website tool, NDC release ver. 1.24.0, 2024. [Online]. Available:
\url{https://seer.cancer.gov/oncologytoolbox/canmed/}. Accessed: Oct. 1, 2024.

\bibitem{seerrx2026}
National Cancer Institute, SEER Program, ``SEER*Rx Interactive Antineoplastic
Drugs Database.'' [Online]. Available:
\url{https://seer.cancer.gov/seertools/seerrx/}. Accessed: Aug. 8, 2026.

\bibitem{shivanna2026uotd}
\ifanonymous
Anonymous, ``Knowledge-base artifact,'' omitted for double-anonymous review,
2026.
\else
A. Shivanna, ``Unified Oncology Treatment Database,'' publication build
\code{publication-build-2026-08-09}, \emph{GitHub}, commit
\code{e4ba3722b5505cfa587b30032b8896d86baf8092}, 2026. [Online]. Available:
\url{https://github.com/shivanna-ornl/unified-oncology-treatment-database/tree/e4ba3722b5505cfa587b30032b8896d86baf8092}.
Accessed: Aug. 11, 2026.
\fi

\bibitem{shivanna2026github}
\ifanonymous
Anonymous, ``Source-code artifact,'' omitted for double-anonymous review, 2026.
\else
A. Shivanna, ``OncoRx-Bench: Oncology drug extraction benchmark,''
\emph{GitHub}, reproducible release, commit
\code{4777e029df40153b422ef7305be4dcf12445958d}, 2026. [Online]. Available:
\url{https://github.com/shivanna-ornl/oncorx-bench/tree/reproducible-v1}.
Accessed: Aug. 11, 2026.
\fi

\bibitem{shivanna2026hf}
\ifanonymous
Anonymous, ``Dataset artifact,'' omitted for double-anonymous review, 2026.
\else
A. Shivanna, ``OncoRx-Bench: Oncology drug extraction benchmark,''
\emph{Hugging Face}, revision
\code{842efd033f21d78b11ee05e6031215e46a440947}, 2026. [Data set].
[Online]. Available:
\url{https://huggingface.co/datasets/AbhishekShivanna/oncorx-bench/tree/842efd033f21d78b11ee05e6031215e46a440947}.
Accessed: Aug. 11, 2026.
\fi

\bibitem{hsu2025}
E. Hsu and K. Roberts, ``LLM-IE: A Python package for biomedical generative
information extraction with large language models,'' \emph{JAMIA Open}, vol. 8,
no. 2, Art. no. ooaf012, 2025, doi:
\href{https://doi.org/10.1093/jamiaopen/ooaf012}{10.1093/jamiaopen/ooaf012}.

\bibitem{brekke2021}
P. H. Brekke, T. Rama, I. Pil\'an, {\O}. Nytr{\o}, and L. {\O}vrelid, ``Synthetic
data for annotation and extraction of family history information from clinical
text,'' \emph{J. Biomed. Semantics}, vol. 12, Art. no. 11, 2021, doi:
\href{https://doi.org/10.1186/s13326-021-00244-2}{10.1186/s13326-021-00244-2}.

\bibitem{li2021}
J. Li, Y. Zhou, X. Jiang, K. Natarajan, S. V. Pakhomov, H. Liu, and H. Xu,
``Are synthetic clinical notes useful for real natural language processing
tasks: A case study on clinical entity recognition,'' \emph{J. Am. Med.
Inform. Assoc.}, vol. 28, no. 10, pp. 2193--2201, 2021, doi:
\href{https://doi.org/10.1093/jamia/ocab112}{10.1093/jamia/ocab112}.

\bibitem{xu2024}
R. Xu \emph{et al.}, ``Knowledge-infused prompting: Assessing and advancing
clinical text data generation with large language models,'' in \emph{Findings
of the Association for Computational Linguistics: ACL 2024}, 2024,
pp. 15496--15523, doi:
\href{https://doi.org/10.18653/v1/2024.findings-acl.916}{10.18653/v1/2024.findings-acl.916}.

\bibitem{singhal2023}
K. Singhal \emph{et al.}, ``Large language models encode clinical knowledge,''
\emph{Nature}, vol. 620, pp. 172--180, 2023, doi:
\href{https://doi.org/10.1038/s41586-023-06291-2}{10.1038/s41586-023-06291-2}.

\bibitem{anthropic2026}
Anthropic, ``Claude,'' 2026. [Online]. Available:
\url{https://www.anthropic.com/claude}. Accessed: Aug. 11, 2026.

\bibitem{openai2026codex}
OpenAI, ``Codex,'' 2026. [Online]. Available:
\url{https://openai.com/codex/}. Accessed: Aug. 11, 2026.

\end{thebibliography}

\end{document}
